\documentclass[11pt,a4paper]{article}

\usepackage[utf8]{inputenc}
\usepackage[T1]{fontenc}
\usepackage{amsmath, amssymb, amsthm}
\usepackage{mathtools}
\usepackage[margin=2.6cm]{geometry}
\usepackage{enumitem}
\usepackage{array}
\usepackage{longtable}
\usepackage[colorlinks=true,linkcolor=blue,citecolor=blue]{hyperref}

\theoremstyle{plain}
\newtheorem{theorem}{Theorem}[section]
\newtheorem{lemma}[theorem]{Lemma}
\newtheorem{corollary}[theorem]{Corollary}
\newtheorem{proposition}[theorem]{Proposition}

\theoremstyle{definition}
\newtheorem{definition}[theorem]{Definition}
\newtheorem{remark}[theorem]{Remark}
\newtheorem*{assertion1}{Assertion 1}
\newtheorem*{assertion2}{Assertion 2}

\newcommand{\R}{\mathbb{R}}
\newcommand{\C}{\mathbb{C}}
\newcommand{\Z}{\mathbb{Z}}
\newcommand{\N}{\mathbb{N}}
\DeclareMathOperator{\ord}{ord}
\DeclareMathOperator{\res}{res}
\DeclareMathOperator{\disc}{disc}
\DeclareMathOperator{\coeff}{coeff}
\DeclareMathOperator{\lc}{lc}
\newcommand{\ideal}[1]{\langle #1\rangle}
\newcommand{\OO}{\mathcal{O}}
\newcommand{\code}[1]{{\ttfamily\small #1}}
\newcommand{\UNCHANGED}{\textbf{[unchanged]}}
\newcommand{\CHANGED}{\textbf{[changed]}}
\newcommand{\NEW}{\textbf{[new]}}
\emergencystretch=1.5em

\title{A Projection Theorem for Cylindrical Algebraic Decomposition\\
via Elimination Ideals\\[1ex]
\large An informal proof extracted from a Lean~4 formalization}
\author{}
\date{June 12, 2026}

\begin{document}
\maketitle

\begin{abstract}
McCallum's projection theorem (Theorem~3.2.3 of his 1985 thesis) is the
mathematical foundation of the McCallum projection operator for cylindrical
algebraic decomposition (CAD): if every coefficient, discriminant and pairwise
resultant arising from a squarefree basis $A$ is order-invariant on a connected
analytic submanifold $S\subseteq\R^{r-1}$, then the elements of $A$ are
analytically delineable on $S$ with pairwise disjoint, order-invariant sections.
We prove a generalization in which the discriminants and resultants are replaced
by \emph{arbitrary nonzero elements of the corresponding elimination ideals}
$\ideal{F,\partial F/\partial x_r}\cap\R[x]$ and $\ideal{F,G}\cap\R[x]$.
The motivation is proof reconstruction: a verifier checking a CAD-based
certificate need not recompute resultants and discriminants---membership of the
projection polynomials in the elimination ideals can be certified by B\'ezout
cofactors and checked by polynomial arithmetic alone.

The proof follows the architecture of McCallum's original argument and reuses
its intermediate results---complexification, Weierstrass preparation, Zariski's
equisingularity theorem (Theorem~4.1.1 of the thesis), and the return to the
real domain---\emph{verbatim}. Exactly three ingredients change:
(i)~the identity $D=\prod_i\disc(F_i)\cdot\prod_{i<j}\res(F_i,F_j)^2$
(thesis Theorem~2.3.3), used to deduce the projection theorem from the lifting
theorem, is replaced by a purely ring-theoretic \emph{radical lemma}: some power
of the product of the chosen elimination-ideal elements again lies in the
elimination ideal of the product;
(ii)~the factorization $E=Q\cdot\disc(h)$ of the discriminant across Weierstrass
preparation (thesis equation~(3.3.8), again via Theorem~2.3.3) is replaced by a
\emph{norm--resultant identity}: any $P$ in $\ideal{h,h'}\cap\OO$ for a
Weierstrass polynomial $h$ of degree $m$ over the ring $\OO$ of analytic germs
satisfies $P^{m}=\pm\disc(h)\cdot Q$ with $Q\in\OO$;
(iii)~the B\'ezout cofactors witnessing $P\in\ideal{f,f'}$ must be transported
through the coordinate change, the complexification and the Weierstrass
preparation, a bookkeeping obligation absent from the original proof.
Everything else, including the entire Chapter~4 of the thesis (the Zariski
theorem), is used unchanged.

The complete proof, including a self-contained proof of Zariski's theorem and
of the Weierstrass preparation and division theorems, has been formalized and
machine-checked in Lean~4 over Mathlib; the main theorem
(\code{mccallum\_3\_2\_3\_generalized}) depends on no axioms beyond Lean's
standard three (\code{propext}, \code{Classical.choice}, \code{Quot.sound}).
This document is the informal rendering of that formal proof.
\end{abstract}

\tableofcontents

%% =========================================================================
\section{Introduction}\label{sec:intro}
%% =========================================================================

McCallum's improved projection operator \cite{McCallum85,McCallum98} replaces
the Collins projection in CAD construction by the much smaller set
\[
  P(A) \;=\; \coeff(A)\,\cup\,\disc(A)\,\cup\,\res(A),
\]
and his Theorem~3.2.3 shows that \emph{order}-invariance of $P(A)$ on a cell
lifts to delineability and order-invariance of $A$ over the cell. The
discriminants and resultants enter the proof through exactly one interface:
$\disc(F)$ (resp.\ $\res(F,G)$) is a nonzero polynomial in the base variables
$x=(x_1,\dots,x_{r-1})$ that vanishes wherever $F,\partial F/\partial x_r$
(resp.\ $F,G$) have a common root---that is, a nonzero element of an
\emph{elimination ideal}. This suggests the generalization proved here: the
theorem survives, with the same conclusions, when each $\disc(F)$ is replaced
by an arbitrary chosen nonzero
$d(F)\in\ideal{F,\partial F/\partial x_r}\cap\R[x]$ and each $\res(F,G)$ by an
arbitrary chosen nonzero $r(F,G)\in\ideal{F,G}\cap\R[x]$.

The practical motivation comes from certificate checking. An untrusted CAD
implementation can emit, along with its answer, the projection polynomials it
used together with B\'ezout cofactors witnessing the memberships
$d(F)=AF+B\,\partial F/\partial x_r$ and $r(F,G)=AF+BG$. A verifier then only
needs to check these polynomial identities---no resultant, subresultant or
discriminant computation, and no commitment to any particular projection
implementation. The generalized theorem is exactly the statement that such
certificates are sound.

Sections~\ref{sec:prelim}--\ref{sec:statements} fix notation and state the two
main results, the generalized lifting theorem (Theorem~\ref{thm:lifting},
``3.2.1$'$'') and the generalized projection theorem
(Theorem~\ref{thm:projection}, ``3.2.3$'$''). Section~\ref{sec:ingredients}
develops the three new ingredients. Sections~\ref{sec:proj-proof}
and~\ref{sec:lifting-proof} give the proofs, following McCallum's \S3.2--3.3
step by step and marking each step \UNCHANGED{} or \CHANGED{} relative to the
thesis. Section~\ref{sec:zariski} recalls the Zariski theorem (thesis
Chapter~4), which is used as a black box with \emph{no} modification, and
records how the formalization proves it. Section~\ref{sec:lean} summarizes the
correspondence between this document and the Lean development.

Throughout, thesis references ``Theorem $m.n.k$'' are to McCallum's 1985 PhD
thesis \cite{McCallum85}.

%% =========================================================================
\section{Preliminaries and notation}\label{sec:prelim}
%% =========================================================================

All background is as in Chapter~2 of the thesis and is used without change.
We single out what is needed, with thesis numbering:
analytic functions and convergent power series (\S2.1); the complexification
of a real-analytic function by its power series (Theorem~2.1.2); the implicit
mapping theorem (Theorem~2.1.7); the Weierstrass preparation theorem
(Theorem~2.1.10); analytic submanifolds and straightening coordinate systems
(Theorem~2.2.1); the behaviour of the vanishing order under analytic
isomorphisms (Theorem~2.3.1); the root continuity principle (Theorem~2.3.2);
and the discriminant-of-a-product formula (Theorem~2.3.3), which we will use
\emph{only} inside the proof of the Zariski theorem, not in the main argument.

Throughout, $x=(x_1,\dots,x_{r-1})$, and a polynomial
$f(x,x_r)\in\R[x,x_r]$ is regarded as univariate in $x_r$ over $\R[x]$. We
write $f'=\partial f/\partial x_r$. For $a\in\R^{r-1}$, $f(a,x_r)\in\R[x_r]$
denotes the specialization.

\begin{definition}[Order; order-invariance (thesis \S2.3, \S3.2)]
For $K=\R$ or $\C$, an analytic function $F$ defined near $p\in K^n$, the
\emph{order} $\ord_p F\in\N\cup\{\infty\}$ is the least $k$ such that some
partial derivative of $F$ of total order $k$ does not vanish at $p$ ($\infty$
if all vanish). $F$ is \emph{order-invariant in} a subset $S$ of its domain if
$p\mapsto\ord_p F$ is constant on $S$. For a polynomial
$g\in\R[x]$ this applies with $F=g$ as a function on $\R^{r-1}$; for
$f\in\R[x,x_r]$, order-invariance ``in a subset of $\R^r$'' refers to $f$ as a
function of all $r$ variables.
\end{definition}

Note that the order at a point of $S$ is always the \emph{ambient} order---the
order of $F$ as a function on a neighbourhood in $K^n$, not of the restriction
$F|_S$ (which is identically zero in the interesting cases).

\begin{definition}[Analytic submanifold; analytic delineability (thesis \S2.2, \S3.2)]
$S\subseteq\R^{r-1}$ is a (nonempty, $s$-dimensional) \emph{analytic
submanifold} if every $p\in S$ has an open neighbourhood $W$ and an analytic
$F:W\to\R^{r-1-s}$, submersive at $p$, with $S\cap W=F^{-1}(0)\cap W$.
A polynomial $f(x,x_r)$ is \emph{analytically delineable} on $S$ if there are
$k\ge 0$ analytic functions $\theta_1<\dots<\theta_k$ on $S$ such that, for
every $a\in S$, the real roots of $f(a,x_r)$ are exactly
$\theta_1(a),\dots,\theta_k(a)$, and the multiplicity of $\theta_i(a)$ as a
root of $f(a,x_r)$ is a constant $m_i\ge 1$ independent of $a$. The graphs of
the $\theta_i$ are the \emph{$f$-sections} over $S$.
\end{definition}

\begin{definition}[Squarefree basis]
A finite set $A\subseteq\R[x,x_r]$ is a \emph{squarefree basis} if every
element has positive degree in $x_r$, is squarefree, and the elements are
pairwise coprime.
\end{definition}

%% =========================================================================
\section{The generalized theorems}\label{sec:statements}
%% =========================================================================

\begin{definition}[Generalized reduced projection]\label{def:proj}
Let $A$ be a squarefree basis. A \emph{generalized reduced projection} of $A$
is any set
\[
  P'(A) \;=\; \coeff(A)\;\cup\;\{\,d(F) : F\in A\,\}\;\cup\;
              \{\,r(F,G) : F,G\in A,\ F\neq G\,\},
\]
where $\coeff(A)$ is the set of nonzero $x_r$-coefficients of elements of $A$,
and $d(F)\in\ideal{F,F'}\cap\R[x]$, $r(F,G)\in\ideal{F,G}\cap\R[x]$ are
\emph{chosen nonzero} elements of the indicated elimination ideals.
\end{definition}

\begin{remark}[Non-vacuity; recovering McCallum's $P(A)$]\label{rem:recover}
The resultant lies in the ideal generated by its arguments (Sylvester matrix),
so $\res(F,F')\in\ideal{F,F'}\cap\R[x]$ and $\res(F,G)\in\ideal{F,G}\cap\R[x]$;
moreover $\res(F,F')=\pm\lc(F)\cdot\disc(F)$. Taking
$d(F)=\res(F,F')$ and $r(F,G)=\res(F,G)$, these are nonzero exactly under the
squarefree-basis hypotheses, and they are order-invariant on $S$ whenever the
elements of McCallum's $P(A)$ are (orders add over products, and
$\lc(F)\in\coeff(A)$). Hence Theorem~\ref{thm:projection} below subsumes thesis
Theorem~3.2.3. Conversely, a nonzero element of $\ideal{F,F'}\cap\R[x]$ exists
\emph{only if} $F$ is squarefree: a repeated factor of $F$ divides every
element of $\ideal{F,F'}$ in $\R(x)[x_r]$ and would force the elimination ideal
to vanish. The membership hypotheses below therefore silently carry the
nondegeneracy (squarefreeness, $\disc\neq 0$) assumptions of the original
statements.
\end{remark}

\begin{theorem}[Generalized lifting theorem, ``3.2.1$'$'']\label{thm:lifting}
Let $f(x,x_r)\in\R[x,x_r]$ and let $S$ be a connected analytic submanifold of
$\R^{r-1}$ on which $f$ is degree-invariant and such that $f(a,x_r)\neq 0$ for
every $a\in S$. Suppose there exists $P\in\R[x]$ with
\[
  P\neq 0,\qquad P\in\ideal{f,\,f'}\cap\R[x],\qquad
  P\ \text{order-invariant in } S .
\]
Then $f$ is analytically delineable on $S$, and for every continuous root
function $\theta$ of $f$ on $S$ (i.e.\ $f(a,\theta(a))=0$ for all $a\in S$),
$f$ is order-invariant in the graph of $\theta$. In particular $f$ is
order-invariant in each $f$-section over $S$.
\end{theorem}

\begin{remark}
Thesis Theorem~3.2.1 is the special case $P=D=\disc(f)$ (more precisely
$P=\res(f,f')$, cf.\ Remark~\ref{rem:recover}); its hypotheses ``$f$ squarefree
with $D\neq 0$'' are subsumed by $P\neq 0$ as explained above. The conclusion
quantified over arbitrary continuous root functions is the form delivered by
the formal proof; it implies (and, given delineability, is equivalent to)
order-invariance in each section.
\end{remark}

\begin{theorem}[Generalized projection theorem, ``3.2.3$'$'']\label{thm:projection}
Let $A$ be a nonempty finite squarefree basis of $r$-variate polynomials,
$r\ge 2$, and let $S$ be a connected analytic submanifold of $\R^{r-1}$.
Suppose each element of $A$ is not identically zero on $S$, and each element
of some generalized reduced projection $P'(A)$ (Definition~\ref{def:proj}) is
order-invariant in $S$. Then each element of $A$ is degree-invariant and
analytically delineable on $S$, the sections of $A$ over $S$ are pairwise
disjoint, and each element of $A$ is order-invariant in every section of $A$
over $S$.
\end{theorem}

%% =========================================================================
\section{The three new ingredients}\label{sec:ingredients}
%% =========================================================================

\subsection{A norm--resultant identity for monic polynomials}\label{ssec:norm}

In McCallum's proof of the lifting theorem, the discriminant $E$ of the
transformed polynomial $g$ is related to the discriminant of its Weierstrass
polynomial $h$ by the explicit factorization $E=Q\cdot\disc(h)$, obtained from
the discriminant-of-a-product formula (thesis Theorem~2.3.3, eq.~(3.3.8)). For
a general witness $P\in\ideal{f,f'}\cap\R[x]$ no such product formula exists.
One might hope that $\ideal{h,h'}\cap R=(\res(h,h'))$ for monic $h$ over a
ring $R$, which would give $\disc(h)\mid P$ directly; this is \textbf{false}:
\begin{equation}\label{eq:counterexample}
  R=\mathbb{Q}[t],\quad h=z^3+t,\quad g=h'=3z^2:\qquad
  3t = 3\cdot h - z\cdot g\in\ideal{h,g}\cap R,\quad
  \res(h,g)=27t^2\nmid 3t .
\end{equation}
The correct statement goes through the algebra norm and costs an $m$-th power.

\begin{lemma}[Norm--resultant identity]\label{lem:norm-res}
Let $R$ be a commutative ring and $h\in R[z]$ monic of degree $m$. For every
$g\in R[z]$, the norm of the multiplication-by-$\bar g$ endomorphism of the
free rank-$m$ $R$-module $R[z]/(h)$ equals the resultant:
\[
  N_{R[z]/(h)\,/\,R}(\bar g)\;=\;\res(h,g).
\]
\end{lemma}

\begin{proof}
Both sides are compatible with base change along ring maps (for the norm,
because the matrix of multiplication by $\bar g$ in the basis
$1,\bar z,\dots,\bar z^{m-1}$ has entries the coefficients of
$g\,z^j \bmod h$, and reduction mod a monic polynomial commutes with ring
maps; for the resultant, by the standard $\res(h^\varphi,g^\varphi)
=\varphi(\res(h,g))$ for monic $h$). By a universal-coefficients reduction it
therefore suffices to prove the identity when $R=K$ is a field over which $h$
splits, $h=\prod_{i=1}^m(z-\alpha_i)$. There
$\res(h,g)=\prod_i g(\alpha_i)$, and the norm factors the same way: by the
multiplicativity of determinants along the filtration induced by
$K[z]/(h)\twoheadrightarrow K[z]/(h/(z-\alpha_1))$, whose kernel is a line on
which $\bar g$ acts by the scalar $g(\alpha_1)$, induction on $m$ gives
$N(\bar g)=\prod_i g(\alpha_i)$.
\end{proof}

\begin{corollary}[Norm identity for elimination ideals]\label{cor:norm-elim}
Let $R$ be a commutative ring, $h\in R[z]$ monic of degree $m$, $g\in R[z]$,
and let $P\in R$ satisfy $P\in\ideal{h,g}\cap R$, say $P=ah+bg$ with
$a,b\in R[z]$. Then
\[
  P^{m} \;=\; \res(h,g)\cdot Q,
  \qquad Q \;=\; N_{R[z]/(h)/R}(\bar b)\;\in R .
\]
In particular, for $g=h'$:\quad $P^{m}=\pm\disc(h)\cdot Q$.
\end{corollary}

\begin{proof}
Reducing $P=ah+bg$ modulo $h$ gives $\bar g\,\bar b=P$ in $R[z]/(h)$. Apply
the norm: it is multiplicative, $N(\bar g)=\res(h,g)$ by
Lemma~\ref{lem:norm-res}, and $N(P)=P^m$ since the scalar $P$ acts as
$P\cdot\mathrm{id}$ on a free module of rank $m$. For $g=h'$ use
$\res(h,h')=\pm\disc(h)$ ($h$ monic).
\end{proof}

Two features of Corollary~\ref{cor:norm-elim} matter in the application.
First, it holds over an \emph{arbitrary} commutative ring; we will apply it
with $R=\OO$ the ring of germs at $0$ of holomorphic functions on
$\C^{r-1}$, so that the cofactor $Q=N(\bar b)$ is automatically a
\emph{holomorphic germ}---being polynomial in the coefficients of $b$ and
$h$---provided the B\'ezout cofactor $b$ has analytic coefficients. This is
why the membership $P\in\ideal{h,h'}$ must be transported with analytic
cofactors (\S\ref{ssec:transport}). Second, the exponent $m$ is harmless: the
order of $P^m$ at a point is $m\cdot\ord P$, so constant order of $P$ along a
set gives constant order of $P^m$.
In the counterexample~\eqref{eq:counterexample}:
$(3t)^3=27t^3=\res(h,g)\cdot t$, with $Q=t$.

\subsection{Order-invariance of factors: Lemma 3.2.2 in ambient form}
\label{ssec:order-add}

The second ingredient is the ``only if'' half of thesis Lemma~3.2.2 (whose
proof the thesis records as ``straightforward''), in precisely the ambient,
along-a-subset form in which the thesis itself applies it, and with the
non-degeneracy hypotheses weakened from ``each factor nonconstant'' to ``each
factor not identically zero''. We include the complete proof, since the
statement is applied in a situation where both factors vanish identically
\emph{on} the subset $S$ (only their ambient orders are finite), and since the
proof is genuinely where connectedness of $S$ enters.

\begin{lemma}[Order-invariance descends to factors]\label{lem:order-add}
Let $U\subseteq\C^n$ be a connected open set, $f,g:U\to\C$ holomorphic and
neither identically zero on $U$. Let $S\subseteq U$ be preconnected, let
$z_0\in S$, and suppose $\ord_z(fg)=\ord_{z_0}(fg)$ for all $z\in S$. Then
$\ord_z f=\ord_{z_0}f$ and $\ord_z g=\ord_{z_0}g$ for all $z\in S$.
\end{lemma}

\begin{proof}
Since $f\not\equiv 0$ on the connected open $U$, the identity theorem gives
$\ord_z f<\infty$ for \emph{every} $z\in U$; likewise for $g$. The order at a
point is a valuation: $\ord_z(fg)=\ord_z f+\ord_z g$. Set $a=\ord_{z_0}f$,
$b=\ord_{z_0}g$, $c=a+b=\ord_{z_0}(fg)$; all are finite, and
$\ord_z f+\ord_z g=c$ for all $z\in S$.

Consider $A=\{z\in S:\ord_z f>a\}$. Because the sum of orders is $c$ and all
orders are finite, $A=\{z\in S:\ord_z g<b\}$. The first description shows $A$
is closed in $S$ ($\ord_z f\ge a+1$ says that all partial derivatives of $f$
of total order $\le a$ vanish at $z$, a closed condition); the second shows
$A$ is open in $S$ ($\ord_z g\le b-1$ says that some particular partial
derivative of order $\le b-1$ is nonzero at $z$, an open condition). Thus $A$
is clopen in the preconnected $S$ and misses $z_0$, hence $A=\emptyset$. By
symmetry $\{z\in S:\ord_z g>b\}=\emptyset$ as well, and since the orders sum
to $c$ everywhere on $S$, $\ord f\equiv a$ and $\ord g\equiv b$ on $S$.
\end{proof}

\begin{remark}\label{rem:ambient}
The hypotheses concern the \emph{ambient} orders at points of $S$ and ambient
non-vanishing ($f\not\equiv0$ on $U$, not on $S$). This is essential: in the
application $f=\disc(h)$ restricted to the section $T^*$ is identically zero
whenever the tracked root has multiplicity $m\ge 2$ (for instance
$f=x_3^2-x_2$ gives a Weierstrass polynomial with $\disc(h)=4x_2$, identically
zero on $T^*=\{x_2=0\}$, of constant ambient order $1$ there). Finiteness of
the ambient orders comes from the identity theorem on $U$, never from
non-vanishing on~$S$.
\end{remark}

\subsection{The radical lemma}\label{ssec:radical}

The third ingredient replaces the use of Theorem~2.3.3
($D=\prod\disc(F_i)\cdot\prod\res(F_i,F_j)^2$) in the deduction of the
projection theorem from the lifting theorem. Its content: the product of the
chosen elimination-ideal elements, raised to a suitable power, is itself an
elimination-ideal element for the product polynomial. The proof is purely
ring-theoretic---no algebraically closed field, no Nullstellensatz, no descent
of coefficients---using only the characterization of the radical of an ideal
as the intersection of the primes containing it.

\begin{lemma}[Radical lemma]\label{lem:radical}
Let $R_0$ be a commutative ring, let $F_1,\dots,F_n\in R_0[x_r]$, and set
$f=\prod_{i=1}^n F_i$. Suppose given
\[
  d_i\in\ideal{F_i,\,F_i'}\cap R_0\quad(1\le i\le n),\qquad
  r_{ij}\in\ideal{F_i,\,F_j}\cap R_0\quad(i\neq j),
\]
and set $\mathcal R=\prod_i d_i\cdot\prod_{i\neq j} r_{ij}\in R_0$. Then
$\mathcal R^{N}\in\ideal{f,\,f'}\cap R_0$ for some $N\ge 1$.
\end{lemma}

\begin{proof}
It suffices to show $\mathcal R\in\sqrt{\ideal{f,f'}}$ in $R_0[x_r]$, since
then $\mathcal R^N\in\ideal{f,f'}$ for some $N$, and $\mathcal R^N\in R_0$.
As $\sqrt{I}=\bigcap\{\mathfrak p\ \text{prime}:\mathfrak p\supseteq I\}$, fix
a prime $\mathfrak p\supseteq\ideal{f,f'}$ of $R_0[x_r]$ and show
$\mathcal R\in\mathfrak p$.

From $f=\prod_i F_i\in\mathfrak p$ and primality, some $F_{i_0}\in\mathfrak p$.
Write $f=F_{i_0}\cdot G$ with $G=\prod_{i\neq i_0}F_i$; then
$f'=F_{i_0}'\,G+F_{i_0}\,G'$, and since $f'\in\mathfrak p$ and
$F_{i_0}G'\in\mathfrak p$, we get $F_{i_0}'\,G\in\mathfrak p$. By primality,
either $F_{i_0}'\in\mathfrak p$, in which case
$\ideal{F_{i_0},F_{i_0}'}\subseteq\mathfrak p$ and hence
$d_{i_0}\in\mathfrak p$; or some $F_{j_0}\in\mathfrak p$ with $j_0\neq i_0$,
in which case $\ideal{F_{i_0},F_{j_0}}\subseteq\mathfrak p$ and hence
$r_{i_0 j_0}\in\mathfrak p$. Either way one factor of $\mathcal R$ lies in
$\mathfrak p$, so $\mathcal R\in\mathfrak p$.
\end{proof}

\begin{remark}
Geometrically, $\mathcal R$ vanishes on the projection of the repeated-root
locus of $f$, which decomposes as
$V(f,f')=\bigcup_iV(F_i,F_i')\cup\bigcup_{i<j}V(F_i,F_j)$; the lemma is the
algebraic shadow of that decomposition. The prime-ideal argument above proves
it directly over $R_0=\R[x]$ (indeed over any commutative ring), avoiding both
the passage to $\overline{K}$ and the descent of the membership back to
$\R[x]$ that a Nullstellensatz argument would require. The exponent $N$ is
immaterial for order-invariance, by additivity of orders.
\end{remark}

%% =========================================================================
\section{Proof of the projection theorem from the lifting theorem}
\label{sec:proj-proof}
%% =========================================================================

This mirrors the thesis proof of Theorem~3.2.3 (\S3.2), with
Lemma~\ref{lem:radical} substituting for Theorem~2.3.3.

\begin{proof}[Proof of Theorem~\ref{thm:projection}]
Degree-invariance of each $F\in A$ on $S$ follows from order-invariance of its
coefficients exactly as in the thesis: the leading coefficient is a nonzero
polynomial of constant order on $S$; if it vanished at some $a\in S$ its order
there would be $\ge 1$, hence $\ge 1$ everywhere on $S$, i.e.\ it would vanish
identically on $S$; applying this argument to successive coefficients and
using that $F$ is not identically zero on $S$ pins the specialized degree.
\UNCHANGED

Disjointness of sections is immediate from pairwise coprimality: a B\'ezout
identity $uF+vG=1$ in $\R(x)[x_r]$, cleared of denominators and specialized at
$a\in S$, prevents common roots of $F(a,x_r)$ and $G(a,x_r)$. \UNCHANGED

Let $f=\prod_{F\in A}F$; it is squarefree (the $F$ are squarefree and pairwise
coprime), of positive degree, degree-invariant on $S$, and $f(a,x_r)\neq 0$
for $a\in S$. \UNCHANGED

\CHANGED\ \emph{The witness.} Let
$\mathcal R=\prod_F d(F)\cdot\prod_{F\neq G}r(F,G)\in\R[x]$. Each factor is
nonzero, so $\mathcal R\neq 0$; each factor is order-invariant in $S$, so
$\mathcal R$ is order-invariant in $S$ (orders add over products). By
Lemma~\ref{lem:radical} there is $N\ge1$ with
$P:=\mathcal R^{N}\in\ideal{f,f'}\cap\R[x]$; and $P\neq0$ is order-invariant
in $S$ since $\ord_a P=N\ord_a\mathcal R$. (In the thesis, this paragraph is
the computation $D=\prod\disc(F_i)\prod\res(F_i,F_j)^2$ followed by
Lemma~3.2.2.)

Applying Theorem~\ref{thm:lifting} to $f$ with witness $P$: $f$ is
analytically delineable on $S$ and order-invariant in the graph of every
continuous root function of $f$ on $S$. Delineability and the
root-multiplicity structure pass from the product $f$ to each factor $F\in A$
as in the thesis (each root of $F(a,\cdot)$ is a root of $f(a,\cdot)$; on
sections, the multiplicities of the coprime factors do not interact); and on
the graph of a continuous root function $\theta$ of any $G\in A$
(a continuous root function of $f$ as well), the order-invariance of the
product $f$ descends to each factor $F$ by Lemma~3.2.2 (in the ambient form,
Lemma~\ref{lem:order-add}, applied to $F$ and its cofactor in $f$, both of
which specialize nontrivially at every point of $S$). \UNCHANGED
\end{proof}

%% =========================================================================
\section{Proof of the lifting theorem}\label{sec:lifting-proof}
%% =========================================================================

We follow the thesis proof of Theorem~3.2.1 (\S3.3) step by step. Recall its
architecture: the theorem reduces to a local statement (Assertion~1), which in
turn reduces to a nonsplitting statement about each individual real root
(Assertion~2); Assertion~2 splits into the case $s=r-1$ ($S$ open) and the
case $1\le s\le r-2$ (positive codimension), and the latter proceeds by
coordinate change, complexification, Weierstrass preparation, the Zariski
theorem, and a return to the real domain.

\subsection{Reduction to the local statements}\label{ssec:reduce}

\UNCHANGED\ The deduction of the theorem from Assertion~1 (local
delineability: each $p\in S$ has a neighbourhood $N$ with $f$ analytically
delineable on $S\cap N$ and order-invariant in each section over $S\cap N$) is
exactly as on thesis pp.~22--23: by connectedness of $S$ the number of
distinct real roots, and the list of their multiplicities, are locally
constant hence constant; the local analytic root functions agree on overlaps
(they enumerate the same roots in increasing order) and patch to global ones;
order-invariance globalizes by a clopen argument on $S$. The reduction of
Assertion~1 to Assertion~2 (nonsplitting of each root cluster, eq.~(3.3.1))
via the root continuity principle (Theorem~2.3.2)---choosing $\epsilon$ below
a third of the minimal root separation of $f(p,x_r)$, circles $C_i$ around the
distinct complex roots $\alpha_i$ of multiplicity $m_i$, and a neighbourhood
$N_0$ on which each $C_i$ encloses exactly $m_i$ roots---is also unchanged, as
is the translation reducing each real root to $\alpha_i=0$. (The translation
$x_r\mapsto x_r+\alpha_i$ fixes $\R[x]$, so it transports the witness $P$ and
its B\'ezout cofactors verbatim.)

For the rest of the section, fix $p\in S$, a real root $\alpha=0$ of
$f(p,x_r)$ of multiplicity $m\ge 1$, and the data $\epsilon$, $C$, $N_0$ as
above. We must produce a neighbourhood $N\ni p$ and an analytic
$\theta:S\cap N\to(-\epsilon,\epsilon)$ with: $f(x,\alpha)=0\iff
\alpha=\theta(x)$ for $x\in S\cap N$, $|\alpha|<\epsilon$; and $f$
order-invariant in the graph of $\theta$.

\subsection{Case $s=r-1$ ($S$ open)}\label{ssec:open}

\CHANGED\ (simplified). Since $S$ is open and connected and $P$ is a nonzero
polynomial of constant order on $S$, $P$ vanishes nowhere on $S$: a nonzero
polynomial cannot vanish identically on a nonempty open set, so its order is
$0$ at some point of $S$, hence at every point. Now write the membership as a
B\'ezout identity
\[
  P(x) \;=\; A(x,x_r)\,f(x,x_r)+B(x,x_r)\,f'(x,x_r)
\]
and specialize at any $a\in S$: since $P(a)\neq 0$, the univariate polynomials
$f(a,x_r)$ and $f'(a,x_r)$ generate the unit ideal, i.e.\ $f(a,x_r)$ is
\emph{separable}. In particular $m=1$ and $f'(p,0)\neq 0$, and the implicit
function theorem (Theorem~2.1.7) yields the analytic simple-root function
$\theta$ near $p$ with $f$ of constant order $1$ on its graph
($f'\neq0$ there). This proves Assertion~2 with constant section order $1$.

(The thesis argument passes through the identity
$(-1)^{m(m-1)/2}R(x)=f_0(x)D(x)$ and the non-vanishing of the leading
coefficient to conclude $R(p)\neq0$; the elimination-ideal hypothesis makes
this detour unnecessary.)

\subsection{Case $1\le s\le r-2$: coordinate change and transport of the
witness}\label{ssec:transport}

\emph{Coordinate change.} \UNCHANGED\ By Theorem~2.2.1 choose a connected
neighbourhood $U\subseteq N_0$ of $p$ and an analytic coordinate system
$\Phi:U\to V$ with
$S\cap U=\{x\in U:\phi_{s+1}(x)=\dots=\phi_{r-1}(x)=0\}$; let
$T=\{y:y_{s+1}=\dots=y_{r-1}=0\}$, so $\Phi(S\cap U)=T\cap V$. Define the
transformed pseudopolynomial
$g(y,y_r)=\sum_j g_j(y)\,y_r^{\,d-j}$, $g_j=f_j\circ\Phi^{-1}$, with analytic
coefficients on $V$.

\CHANGED\ \emph{Transport of the witness and of its certificate.} Set
$\widetilde P=P\circ\Phi^{-1}$, analytic on $V$ and not identically zero.
Substituting $x=\Phi^{-1}(y)$ in the B\'ezout identity gives
\begin{equation}\label{eq:transported-membership}
  \widetilde P(y)\;=\;\widetilde A(y,y_r)\,g(y,y_r)
    +\widetilde B(y,y_r)\,\frac{\partial g}{\partial y_r}(y,y_r),
\end{equation}
where $\widetilde A,\widetilde B$ are polynomials in $y_r$ with coefficients
analytic in $y$ (the substitution touches only the base variables, so it
commutes with $\partial/\partial x_r$). Thus
$\widetilde P\in\ideal{g,\partial g/\partial y_r}$ over the ring of analytic
functions---\emph{with explicit analytic cofactors}, which the sequel must
preserve. By Theorem~2.3.1 (order is preserved by analytic isomorphisms),
$\widetilde P$ is order-invariant in $T\cap V$, with the same constant order
$\mu:=\ord_0\widetilde P<\infty$ as $P$ on $S$ near $p$.
(In the thesis this step defines $E=\disc(g)=D\circ\Phi^{-1}$ instead; no
cofactor transport is needed there.)

\subsection{Complexification}\label{ssec:complexify}

\UNCHANGED\ in method (applied to $\widetilde P,\widetilde A,\widetilde B$ in
place of $E$). Expanding all the analytic coefficient functions
($g_j$, $\widetilde P$, and the $y$-coefficients of $\widetilde A,\widetilde
B$) in convergent power series at $0$ and substituting complex variables
(Theorem~2.1.2) extends everything to a polydisc
$\Delta_1\subseteq\C^{r-1}$: a pseudopolynomial $g(z,z_r)$, a holomorphic
$\widetilde P(z)$, and identity~\eqref{eq:transported-membership} continues to
hold on $\Delta_1$ by the identity theorem (both sides are holomorphic and
agree on the real points, coefficientwise in $y_r$).

Let $T^*=\{z:z_{s+1}=\dots=z_{r-1}=0\}$. The thesis argument on pp.~25--26
applies verbatim to $\widetilde P$: every partial derivative of
$\widetilde P$ of order $<\mu$ vanishes on $T\cap V$ by order-invariance, its
complexification vanishes on $T^*\cap\Delta_1$ (a power series with real
coefficients vanishing for real arguments vanishes identically), and some
order-$\mu$ derivative is nonzero at $0$, hence on a polydisc. After
refining $\Delta_1$:
\begin{equation}\label{eq:Ptilde-oi}
  \widetilde P(z)\ \text{is order-invariant in}\ T^*\cap\Delta_1,
  \ \text{with constant (ambient) order}\ \mu .
\end{equation}
Moreover the order of the complexification at a real point equals the real
order (the complex Taylor coefficients at a real point are determined by the
real ones), which is what makes~\eqref{eq:Ptilde-oi} usable and what will
later transport the complex conclusions back to $\R$.

\subsection{Weierstrass preparation and descent of the membership}
\label{ssec:prep}

\emph{Preparation.} \UNCHANGED\ Since $z_r=0$ is a zero of $g(0,z_r)$ of
multiplicity exactly $m$ (and $g(0,z_r)\neq0$ for $0<|z_r|\le\epsilon$), the
Weierstrass preparation theorem (Theorem~2.1.10) gives, after refining to
$\Delta_2\subseteq\Delta_1$, a non-vanishing holomorphic unit $u(z,z_r)$ on
$\Delta'=\Delta_2\times\Delta(0;\epsilon)$ and a Weierstrass polynomial
\[
  h(z,z_r)=z_r^{m}+a_1(z)z_r^{m-1}+\dots+a_m(z),\qquad a_i(0)=0,
\]
with $g=u\,h$ on $\Delta'$ (thesis eq.~(3.3.7)), all roots of $h(z,\cdot)$
lying in $\Delta(0;\epsilon)$.

\CHANGED\ \emph{Descent of the membership to $\ideal{h,h'}$ over germs.} Let
$\OO$ denote the ring of germs at $0\in\C^{r-1}$ of holomorphic functions; it
is an integral domain (identity theorem), and $h\in\OO[z_r]$ is monic of
degree $m$. We claim
\begin{equation}\label{eq:descent}
  \widetilde P\;\in\;\ideal{h,\,h'}\ \subseteq\ \OO[z_r]
  \qquad(h'=\partial h/\partial z_r).
\end{equation}
Indeed, differentiating $g=uh$ in $z_r$ gives $g'=u'h+uh'$;
substituting into~\eqref{eq:transported-membership},
\[
  \widetilde P
  \;=\;\bigl(\widetilde A\,u+\widetilde B\,u'\bigr)\,h
      +\bigl(\widetilde B\,u\bigr)\,h' ,
\]
an identity of holomorphic functions of $(z,z_r)$ near $0$ whose cofactors
$A_1=\widetilde A u+\widetilde B u'$ and $B_1=\widetilde B u$ are holomorphic
but not polynomial in $z_r$. To restore polynomiality, divide $B_1$ by the
monic $h$ (Weierstrass division): $B_1=qh+b$ with $b\in\OO[z_r]$ of degree
$<m$ and $q$ holomorphic. Then
$\widetilde P-b\,h'=(A_1+q\,h')\,h$, and the left side is a polynomial in
$z_r$ of degree $\le 2m-2$ with germ coefficients; dividing by the monic $h$
in $\OO[z_r]$ and using that $\OO$ is a domain, the holomorphic quotient
$A_1+q h'$ coincides with the polynomial quotient $a\in\OO[z_r]$. Hence
$\widetilde P=a\,h+b\,h'$ with $a,b\in\OO[z_r]$,
proving~\eqref{eq:descent}---with the crucial byproduct that the cofactor $b$
has \emph{analytic} coefficients.

\subsection{The discriminant of $h$ has constant order along $T^*$}
\label{ssec:disc-order}

\CHANGED\ (this is the key new step, replacing thesis eq.~(3.3.8)).
Apply Corollary~\ref{cor:norm-elim} over $R=\OO$ to the
membership~\eqref{eq:descent}:
\begin{equation}\label{eq:norm-identity}
  \widetilde P(z)^{\,m}\;=\;\res(h,h')(z)\cdot Q(z)
  \;=\;\pm\,\disc(h)(z)\cdot Q(z),
  \qquad Q=N_{\OO[z_r]/(h)}(\bar b)\in\OO ,
\end{equation}
valid on a connected open neighbourhood $\Delta_3\subseteq\Delta_2$ of $0$ on
which all the germs involved have analytic representatives. The cofactor $Q$
is analytic because the norm is a polynomial expression in the coefficients of
$b$ and of $h$ (\S\ref{ssec:norm}).

Neither factor on the right of~\eqref{eq:norm-identity} vanishes identically
on $\Delta_3$: since $\widetilde P\not\equiv 0$ (it has finite order $\mu$ at
$0$) we have $\widetilde P^m\not\equiv0$, and $\OO$ is a domain, so
$\disc(h)\not\equiv0$ and $Q\not\equiv0$. (Note that squarefreeness of $f$ is
never invoked: the non-vanishing of $\disc(h)$ is a consequence of the
existence of the nonzero witness, cf.\ Remark~\ref{rem:recover}.)

By~\eqref{eq:Ptilde-oi} and additivity of orders,
$\ord_z\widetilde P^{\,m}=m\mu$ is constant for $z\in T^*\cap\Delta_3$. Now
apply Lemma~\ref{lem:order-add} on $U=\Delta_3$ with $S=T^*\cap\Delta_3$
(preconnected), $f=\disc(h)$, $g=Q$:
\begin{equation}\label{eq:disc-oi}
  \disc(h)\ \text{is order-invariant in}\ T^*\cap\Delta_3,
  \ \text{with finite constant order} .
\end{equation}
By Remark~\ref{rem:ambient} the lemma applies even though
$\disc(h)|_{T^*}\equiv 0$ whenever $m\ge2$---the generic situation.
Display~\eqref{eq:disc-oi} is precisely the hypothesis of Zariski's theorem.

(For comparison, the thesis obtains~\eqref{eq:disc-oi} from
$E=Q\cdot\disc(h)$ with $Q=\disc(u)\res(u,h)^2$ via Theorem~2.3.3, plus
Lemma~3.2.2; here the unit $u$ never needs to be analyzed as a
pseudopolynomial, and the witness replaces the discriminant $E$ throughout.)

\subsection{Application of Zariski's theorem}\label{ssec:zariski-apply}

\UNCHANGED\ By~\eqref{eq:disc-oi}, Theorem~4.1.1 of the thesis (stated as
Theorem~\ref{thm:zariski} below) applies to the Weierstrass polynomial $h$:
after refining to $\Delta_4\subseteq\Delta_3$ there is a holomorphic
$\psi:\Delta_4^{(s)}\to\Delta(0;\epsilon)$ such that for
$z=(z_1,\dots,z_s,0,\dots,0)\in T^*\cap\Delta_4$ and
$\alpha\in\Delta(0;\epsilon)$,
\[
  h(z,\alpha)=0\iff\alpha=\psi(z_1,\dots,z_s)
  \qquad\text{(thesis eq.~(3.3.9)),}
\]
the unique root having multiplicity exactly $m$, and $h$ is order-invariant in
the graph $G^*=\{(z,\psi(z)):z\in T^*\cap\Delta_4\}$. This delivers both the
\emph{nonsplitting} of the root cluster over the section and the
order-invariance of $h$ (hence of $g=uh$, as $u\neq0$) along it.

\subsection{Return to the real domain and reassembly}\label{ssec:real}

\UNCHANGED\ The argument of thesis pp.~27--28 applies verbatim:
\begin{enumerate}[leftmargin=2em]
\item For real $y\in T\cap B$ the polynomial $g(y,y_r)$ has real
coefficients, so the unique root $\psi(y)$ of $h(y,\cdot)$ in
$\Delta(0;\epsilon)$ is fixed by conjugation, i.e.\ real. Consequently the
power series of $\psi$ has real coefficients and the restriction
$\eta=\psi|_{B^{(s)}}$ is a real-valued real-analytic function (Schwarz
reflection, via Theorem~2.1.2).
\item For $y\in T\cap B$, $\alpha\in(-\epsilon,\epsilon)$:
$g(y,\alpha)=0\iff h(y,\alpha)=0\iff\alpha=\eta(y)$, since $u\neq0$.
\item The order of $g$ at a real point of the graph of $\eta$ equals the
order of its complexification at that point (matching Taylor coefficients,
as in \S\ref{ssec:complexify}); hence the order-invariance of $h$ in $G^*$
gives order-invariance of $g$ in the graph of $\eta$.
\item Detransform: $N=\Phi^{-1}(B)$ and $\theta=\eta\circ\phi$, where $\phi$
is the chart of $S$ induced by $\Phi$. By Theorem~2.3.1 the order is
preserved, establishing Assertion~2 for the root $\alpha_i$.
\end{enumerate}
Undoing the translation of \S\ref{ssec:reduce} and intersecting the finitely
many neighbourhoods $N_i$ yields Assertion~1 at $p$; globalization over the
connected $S$ was already described in \S\ref{ssec:reduce}. Finally, for the
conclusion as stated (order-invariance along \emph{arbitrary} continuous root
functions $\theta$): on each connected local piece, the values of $\theta$ lie
among the finitely many separated analytic branches, so the branch index is
locally constant, hence constant; thus $\theta$ coincides with a single
delineation branch and inherits its order-invariance. \qed

%% =========================================================================
\section{The Zariski theorem (thesis Chapter 4): unchanged}
\label{sec:zariski}
%% =========================================================================

The generalization makes \emph{no} change to Chapter~4 of the thesis: the
hypothesis it consumes, constancy of the order of $\disc(h)$ along the
section, is produced by \S\ref{ssec:disc-order}, and its conclusion is
consumed in \S\ref{ssec:zariski-apply} exactly as in the original proof. We
restate it for completeness.

\begin{theorem}[Zariski; thesis Theorem 4.1.1]\label{thm:zariski}
Let $h(z,z_n)$ be a Weierstrass polynomial of degree $m\ge 1$ on a polydisc
$\Delta_1\subseteq\C^{n-1}$, all of whose roots lie in $\Delta(0;\epsilon)$,
and whose discriminant $F(z)$ does not vanish identically. Let
$1\le s\le n-2$ and $T^*=\{z:z_{s+1}=\dots=z_{n-1}=0\}$, and assume $F$ is
order-invariant in $T^*\cap\Delta_1$. Then on a smaller polydisc there is a
holomorphic $\psi(z_1,\dots,z_s)$ with values in $\Delta(0;\epsilon)$ such
that, over points of $T^*$, the roots of $h$ are exactly $z_n=\psi$
(``nonsplitting'': one root of multiplicity $m$), and $h$ is order-invariant
in the graph $G^*$ of $\psi$ over $T^*$.
\end{theorem}

Since the formalization had to be axiom-free, the entire Chapter-4 proof was
formalized as well, following the thesis: the splitting of Theorem~4.2.1 into
the nonsplitting statement (Theorem~4.2.2) and the order-invariance statement
(Theorem~4.2.3); factorization of $h$ into irreducible Weierstrass polynomials
over the germ ring; the root cover of the separable locus and path lifting
(Lemma~4.2.4); transitive monodromy for irreducible factors (Lemma~4.2.5);
the homotopy-deformation contradiction proving nonsplitting per factor and the
resultant/zero-set argument making the factors' roots coincide
(proof of Theorem~4.2.2); the Newton--Puiseux parametrization
$z_{n-1}=u^m$, $z_n=\phi(z,u)$ (Lemma~4.2.6) and the order computation
$\ord_{(z,0,0)}h=\min(m,m_1)$ (Lemmas~4.2.7--4.2.8, proof of Theorem~4.2.3);
and, for codimension $\ge 2$, the quadratic blow-up
$Q(z,z_{s+1},\dots,z_{n-1})=(z,\,z_{s+1}z_{n-1},\dots,z_{n-2}z_{n-1},\,
z_{n-1})$ reducing to codimension one, including the normal form
$F\circ Q=z_{n-1}^{\,r}\,N$ with $N(0)\neq0$ (stages 1--7 of the thesis proof
of 4.1.1) and the term-tracking transport of the order of $h$ through the
blow-up (stage 8). A few proofs were reorganized in the formalization without
changing the statements; the noteworthy points are:

\begin{itemize}[leftmargin=2em]
\item \emph{Lemma 4.2.5.} The thesis cites Bochner--Martin for the transitive
monodromy of the root cover of an irreducible Weierstrass polynomial. The
formalization proves it: a clopen splitting of the root cover over the
punctured (separable) base would glue, via Cauchy-coefficient continuity and a
removable-singularity argument, into a nontrivial monic factorization of $h$
over the germ ring, contradicting irreducibility; path-connectedness of the
cover then yields transitive monodromy by lifting.
\item \emph{Genericity of the blow-up direction} (stage 1 of the thesis
proof): the linear change making $\partial^r F/\partial z_{n-1}^r(0)\neq 0$ is
obtained from a polarization argument---the degree-$r$ leading form of $F$ is
a nonzero symmetric form, so its diagonal restriction cannot vanish on any
dense set of transverse directions. The same mechanism drives the stage-8
term-tracking (the nonzero coefficient $q$ of the thesis's displayed term is
realized at some nearby ratio point $w''$).
\item \emph{Lemmas 4.2.7--4.2.8.} Where the thesis works in iterated
fractional-power-series fields, the formalization stays with the convergent
parametrization of Lemma~4.2.6 and uses only \emph{continuity} of
Cauchy-integral coefficients in the section variables, which suffices for the
order computation $\ord_{(z,0,0)}h=\min(m,m_1)$.
\item The auxiliary identities of Theorem~2.3.3 (discriminant of a product)
\emph{are} used here---inside Chapter~4, to descend discriminant
order-invariance from $h$ to its irreducible factors and to force distinct
factors to share their section root---exactly as in the thesis proofs of
Theorems~4.2.2 and~4.2.3.
\end{itemize}

We emphasize the division of labour this creates for the generalized theorem:
all uses of the \emph{specific} polynomials $\disc$, $\res$ in McCallum's
Chapter~3 argument are eliminated in favour of the elimination-ideal witness,
while their uses \emph{internal} to the Zariski theorem are untouched---they
concern the Weierstrass polynomial $h$ and its factors, objects private to the
proof, not the projection set handed to the CAD algorithm.

%% =========================================================================
\section{Summary of changes, and the Lean formalization}\label{sec:lean}
%% =========================================================================

\subsection{What changed relative to the thesis}

\begin{center}
\renewcommand{\arraystretch}{1.25}
\begin{tabular}{|p{5.4cm}|c|p{7.2cm}|}
\hline
\textbf{Thesis item} & \textbf{Status} & \textbf{Notes} \\
\hline
Ch.~2 (preliminaries) & unchanged & used as is \\
\hline
Def.\ of $P(A)$ & changed & $P'(A)$ with chosen elimination-ideal elements
(Def.~\ref{def:proj}) \\
\hline
Thm.~3.2.1 statement & changed & witness $P\in\ideal{f,f'}\cap\R[x]$ replaces
$D$; squarefree/$D\neq0$ subsumed by $P\neq 0$ \\
\hline
Thm.~3.2.3 statement & changed & $P'(A)$ replaces $P(A)$ \\
\hline
Proof of 3.2.3 from 3.2.1 & changed & radical lemma (Lemma~\ref{lem:radical})
replaces $D=\prod\disc\cdot\prod\res^2$ (Thm.~2.3.3) \\
\hline
3.2.1: reduction to Assertions 1, 2; root clusters; translation &
unchanged & \S\ref{ssec:reduce} \\
\hline
3.2.1: case $s=r-1$ & simplified & $P(a)\neq0$ $\Rightarrow$ separability
directly (\S\ref{ssec:open}) \\
\hline
3.2.1: coordinate change & unchanged & \S\ref{ssec:transport} \\
\hline
3.2.1: transform of the invariant & changed & $\widetilde P=P\circ\Phi^{-1}$
\emph{with B\'ezout cofactors} replaces $E=\disc g$ \\
\hline
3.2.1: complexification & unchanged in method & applied to $\widetilde P$
(and cofactors) instead of $E$ \\
\hline
3.2.1: Weierstrass preparation & unchanged & eq.~(3.3.7) \\
\hline
3.2.1: membership descent to $\ideal{h,h'}$ & new & Weierstrass division of
the cofactors (\S\ref{ssec:prep}) \\
\hline
3.2.1: factorization step (3.3.8) & changed & norm identity
$\widetilde P^{m}=\pm\disc(h)\,Q$ (Cor.~\ref{cor:norm-elim}) $+$
Lemma~\ref{lem:order-add} replace Thm.~2.3.3 \\
\hline
3.2.1: Zariski application (3.3.9) & unchanged & \S\ref{ssec:zariski-apply} \\
\hline
3.2.1: return to real; reassembly & unchanged & \S\ref{ssec:real} \\
\hline
Ch.~4 (Zariski theorem) & unchanged & statement and proof as in thesis;
formalized in full (\S\ref{sec:zariski}) \\
\hline
\NEW\ norm--resultant identity & new & Lemma~\ref{lem:norm-res},
Cor.~\ref{cor:norm-elim} (classical algebra) \\
\hline
\NEW\ order-additivity lemma & new & Lemma~\ref{lem:order-add} = ambient
``only if'' half of Lemma~3.2.2, full proof \\
\hline
\NEW\ radical lemma & new & Lemma~\ref{lem:radical} (pure ring theory) \\
\hline
\end{tabular}
\end{center}

\subsection{The formalization}

The entire development is machine-checked in Lean~4 over Mathlib. The main
theorem is \code{mccallum\_3\_2\_3\_generalized}
(\code{Mccalum/Generalized/Projection.lean}); \code{\#print axioms} reports
only \code{propext}, \code{Classical.choice}, \code{Quot.sound}. In
particular, the classical inputs that the informal literature cites---the
Weierstrass preparation and division theorems for convergent power series
(absent from Mathlib), the several-variable
holomorphic$\Rightarrow$analytic bridge needed for them, Zariski's theorem
with its monodromy and blow-up arguments---are all proved inside the
development. A correspondence table for the main steps:

\begin{center}
\renewcommand{\arraystretch}{1.2}
\small
\begin{tabular}{|p{6.0cm}|p{8.4cm}|}
\hline
\textbf{This document} & \textbf{Lean identifier (file)} \\
\hline
Theorem~\ref{thm:projection} (3.2.3$'$) &
\code{mccallum\_3\_2\_3\_generalized} (\code{Generalized/Projection.lean}) \\
\hline
Theorem~\ref{thm:lifting} (3.2.1$'$) &
\code{lifting\_theorem\_generalized} (ibid.) \\
\hline
Lemma~\ref{lem:radical} &
\code{nullstellensatz\_elim} (ibid.) \\
\hline
Lemma~\ref{lem:norm-res} &
\code{norm\_eq\_resultant\_monic} (\code{Generalized/Lifting.lean}) \\
\hline
Corollary~\ref{cor:norm-elim} &
\code{norm\_identity\_elim} (ibid.) \\
\hline
Lemma~\ref{lem:order-add} &
\code{order\_factor\_const\_of\_mul\_analytic} (\code{OrderMulAnalytic.lean}) \\
\hline
\S\ref{ssec:open} (open case) &
\code{lifting\_generalized\_open\_case},
\code{separable\_of\_elim\_nonvanishing} (\code{Generalized/Lifting.lean}) \\
\hline
\S\ref{ssec:transport} (chart + transport) &
\code{lifting\_generalized\_codim\_local} (\code{Generalized/Delineable.lean}) \\
\hline
\S\ref{ssec:complexify} &
\code{analyticAt\_complexify}, \code{order\_real\_eq\_order\_complex}
(\code{Generalized/Lifting.lean}) \\
\hline
\S\ref{ssec:prep} (preparation; descent) &
\code{weierstrass\_preparation\_analytic} (\code{Generalized/WeierstrassPrep.lean});
\code{descent\_membership} (\code{Generalized/MembershipDescent.lean}) \\
\hline
Eq.~\eqref{eq:norm-identity} over germs &
\code{descent\_norm\_identity} (\code{Generalized/DescentNormIdentity.lean}) \\
\hline
\S\ref{ssec:disc-order} (disc order constant) &
\code{weierstrassDisc\_order\_const\_along\_section}
(\code{Generalized/DiscOrder.lean}) \\
\hline
Theorem~\ref{thm:zariski} as consumed &
\code{cluster\_root\_structure} (\code{Generalized/ClusterRootStructure.lean})
= \code{zariski\_single\_branch} (Conclusion~1,
\code{Generalized/ZariskiNonsplitting.lean}) $+$
\code{Puiseux.order\_invariant\_in\_graph} (Conclusion~2,
\code{Puiseux/Conclusion2General.lean}) \\
\hline
Thesis Thm.~4.2.2 (nonsplitting, codim 1) &
\code{irreducible\_section\_single\_root\_e1} (\code{Generalized/ZariskiE1.lean});
monodromy kernel in \code{Generalized/Monodromy*.lean},
\code{Puiseux/RootCover.lean} \\
\hline
Thesis Lemma 4.2.5 &
\code{rootCover\_exchange\_gen} (\code{Generalized/MonodromyGen.lean}) \\
\hline
Thesis Thm.~4.2.3 / Lemmas 4.2.6--4.2.8 &
\code{Puiseux/Parametrization.lean}, \code{Puiseux/OrderInvariance.lean},
\code{Puiseux/Conclusion2Value*.lean} (\code{order\_eval\_eq\_min}) \\
\hline
Thesis 4.1.1 Case II (blow-up; stage 8) &
\code{ZariskiE2.irreducible\_section\_single\_root\_blowup\_proof}
(\code{Generalized/ZariskiE2Main.lean});
\code{Puiseux.order\_eval\_value\_e2\_blowup\_proof}
(\code{Puiseux/Conclusion2E2.lean}) \\
\hline
\S\ref{ssec:real}, reassembly, globalization &
\code{Generalized/ClusterFromReal.lean}, \code{MultiCluster.lean};
\code{locally\_delineable\_to\_global$'$},
\code{order\_invariant\_of\_locally\_invariant}
(\code{Generalized/Delineable.lean}) \\
\hline
\end{tabular}
\end{center}

Two findings of the formalization are worth recording. First, the natural
strengthening ``$\ideal{h,g}\cap R=(\res(h,g))$ for monic $h$'' that one might
be tempted to use in \S\ref{ssec:disc-order} is false
(counterexample~\eqref{eq:counterexample}); the formal proof effort exposed
this early, and Lemma~\ref{lem:norm-res} is the correct replacement. Second,
the deduction of the projection theorem needs no Nullstellensatz and no
algebraically closed field: the radical-as-intersection-of-primes argument of
Lemma~\ref{lem:radical} is both simpler to formalize and sharper than the
geometric argument it replaces.

\begin{thebibliography}{9}
\bibitem{McCallum85} S.~McCallum,
  \emph{An Improved Projection Operation for Cylindrical Algebraic
  Decomposition}, PhD thesis, University of Wisconsin--Madison, 1985.
\bibitem{McCallum98} S.~McCallum,
  \emph{An improved projection operation for cylindrical algebraic
  decomposition}, in: B.~Caviness, J.~Johnson (eds.), \emph{Quantifier
  Elimination and Cylindrical Algebraic Decomposition}, Springer, 1998,
  242--268.
\bibitem{Zariski65} O.~Zariski, \emph{Studies in equisingularity I--II},
  Amer.\ J.\ Math.\ \textbf{87} (1965).
\bibitem{Zariski75} O.~Zariski,
  \emph{On equimultiple subvarieties of algebroid hypersurfaces},
  Proc.\ Nat.\ Acad.\ Sci.\ USA \textbf{72} (1975), 1425--1426.
\bibitem{Lang} S.~Lang, \emph{Algebra}, 3rd ed., Springer GTM 211, Ch.~IV
  (resultants and norms).
\end{thebibliography}

\end{document}
